% =============================================================================
%  06_derived.tex —— §6 导出函子
% =============================================================================
\section{导出函子}\label{sec:df}

现在可以回到（导出）函子：把任意复形换成适合所论函子的对象组成的消解。我们处理
左正合函子；右正合函子类似。

\begin{definition}\label{def:adapted}
设 $\A$、$\B$ 是阿贝尔范畴。对象类 $\RR\subset\A$ 称为\term{适应}{adapted}于
左正合函子 $F\colon\A\to\B$（\cite{GM} III 6.3），如果
\begin{itemize}
\item $\RR$ 对直和封闭；
\item 把 $F$ 作用到 $\RR$ 中的无环复形（即上同调为零的复形）上，所得复形仍无环；
\item 任何 $A\in\A$ 都能通过 $0\to A\to R$ 嵌入某个 $R\in\RR$。
\end{itemize}
再令 $K^+(\RR)$ 为 $\RR$ 中有界下方链复形构成的范畴，态射取链映射的同伦类。
则把 $K^+(\RR)$ 中的拟同构可逆化（\cite{GM}），所得范畴等价于 $D^+(\A)$。
\end{definition}

最后这句话说得花哨些，意思无非是：目标终于达到了——利用定义中的条件，我们可以
函子化地把任何 $\A$-复形换成拟同构的 $\RR$-复形（即作消解）。这样的
$\RR\subset\A$ 的例子有：$\{$内射层$\}\subset\{$拟凝聚层$\}$；或者对右正合
函子，有 $\{$环上的投射模$\}\subset\{$全体模$\}$、
$\{$仿射簇或仿射概形上的局部自由层$\}\subset\{$凝聚层$\}$、
$\{$平坦层$\}\subset\{$拟凝聚层$\}$（当 $F$ 是与某个层作张量积时）等等。

于是我们不再把 $F$ 作用到任意复形上，而只作用到 $\RR$-复形上。

这样便定义 $F$ 的\term{右导出函子}{right derived functor} $\R F$ 为复合
$$
D^+(\A)\longrightarrow K^+(\RR)/\text{拟同构}
\xrightarrow{\ F\ }D^+(\B).
$$
（更经典的右导出函子 $R^iF$ 就是上同调 $R^iF=h^i(\R F)$；不过我们现在知道，
取上同调是坏事。）这里略过了一些细节：严格说来，$\R F$ 应当用某个泛性质来定义
（\cite{GM} III 6.11），才能保证它与 $\RR$ 的选取无关，只差典范同构。同样可以
对右正合函子 $F$ 定义\term{左导出函子}{left derived functor} $\mathbf{L}F$。

这给出一个\emph{正合函子} $\R F\colon D^+(\A)\to D^+(\B)$，即它把正合三角送到
正合三角。特别地，取上同调就得到经典的导出函子长正合列
$$
\ldots\to R^iF(A)\to R^iF(B)\to R^iF(C)\to R^{i+1}F(A)\to\ldots
$$

例如，左正合的整体截面函子
$\Gamma\colon\{\text{层}\}\to\{\text{向量空间}\}$ 的右导出函子就是层上同调
$\R\Gamma$，它的（向量空间复形的）上同调就是标准的层上同调 $H^i$。于是上面的
序列就成为上同调的长正合列。

这种做法有两大好处。第一，我们把\emph{复形} $\R F(A)$——而不是它那较弱的
上同调 $R^iF(A)$——变成了 $A$ 的不变量，且在拟同构意义下唯一。第二，导出函子
就是原来的函子作用到复形上（当然不是任意复形，必须取自 $\RR$）。这就为导出
函子的结论给出了更简单、更有概念性的证明，而传统上这些结论需要复杂的双复形、
谱序列论证。下面在层论中举几个例子，其间略过一些技术条件（关于复形与消解有界性
的问题，在光滑射影簇上当然不成问题；精确陈述见 \cite{HRD}）：

\begin{eg}\label{eg:tor}
层的张量积是对称的，因此它的导出函子 $\Ltimes$ 与同调 $\Tor_i$ 也是对称的：
$\Tor_i(A,B)\cong\Tor_i(B,A)$。这里我们只是把平坦层（例如局部自由层）组成的
复形作张量积。
\end{eg}

\begin{eg}\label{eg:rhom}
层的 $\R\mathcal Hom$ 无非是复形上的局部 $\mathcal Hom$，可以用局部自由层消解
第一个变量，或用内射对象消解第二个变量来定义。
\end{eg}

\begin{eg}\label{eg:composition}
在一些温和条件下，$\R(F\circ G)\cong\R F\circ\R G$。若在作用 $\R F$ 之前先取
上同调，就得到 $\R(F\circ G)$ 的一个近似，即 Grothendieck 谱序列
$$
R^iF(R^jG)\Rightarrow R^{i+j}(F\circ G).
$$

例如对态射 $p\colon X\to Y$，等式 $\Gamma_X=\Gamma_Y\circ p_*$ 给出
$\R\Gamma_X\cong\R\Gamma_Y\circ\R p_*$，于是有 Leray 谱序列。又
$\Hom=\Gamma\circ\mathcal Hom$ 给出
$\R\Hom\cong\R H\udot(\R\mathcal Hom)$ 与局部到整体的谱序列
$\Ext^{i+j}\Leftarrow H^i(\mathcal Ext^j)$。

超上同调 $\mathbb H$ 由“复形的整体截面”这个导出函子给出（对我们来说，它无非
是通常的层上同调）。把上面的讨论用于 $F=\Gamma$、$G=\id$，就得到超上同调谱序列
$H^i(h^j(A\udot))\Rightarrow\mathbb H^{i+j}(A\udot)$。
\end{eg}

\begin{eg}\label{eg:duality}
设 $A$ 是层或层的复形，记 $A\ve$ 为对偶复形
$\R\mathcal Hom(A,\OO)$。则 $\R\mathcal Hom(A,B)\cong B\Ltimes A\ve$。
\end{eg}

\begin{eg}\label{eg:derham}
现在看实 $n$ 维流形上的一个例子，即 de Rham 定理。设 $A^i(\mathbb R)$ 表示流形
$X$ 上 $C^\infty$ $i$-形式组成的层，$d$ 为外微分。Poincaré 引理给出拟同构
（常数层 $\mathbb R$ 的消解）
$$
\mathbb R\to\{A^0(\mathbb R)\xrightarrow{d}A^1(\mathbb R)\xrightarrow{d}\ldots
\xrightarrow{d}A^n(\mathbb R)\}.
$$
由于单位分解的存在，层 $A^i(\mathbb R)$ 对整体截面函子 $\Gamma$ 无环，即它们
没有高阶层上同调（这类层称为\term{细层}{fine sheaf}）。因此容易看出，可以把
$\Gamma$ 作用到这个消解上，得到与常数层 $\mathbb R$ 的导出函子
$\R\Gamma(\mathbb R)$（即它的层上同调）同构的复形。这就给出
$$
\Omega^0(\mathbb R)\xrightarrow{d}\Omega^1(\mathbb R)\xrightarrow{d}\ldots
\xrightarrow{d}\Omega^n(\mathbb R),
$$
也就是说，de Rham 复形计算了流形 $X$ 的实系数上同调 $H^*(X;\mathbb R)$。

同样，在复流形 $X$ 上，对下述复形（原文称之为 Dolbeault 复形或消解，见校勘
\ref{kanwu:dolbeault}）取上同调，
$$
\C\to\{\OO\xrightarrow{\partial}\Omega^{1,0}\xrightarrow{\partial}\Omega^{2,0}
\xrightarrow{\partial}\ldots\xrightarrow{\partial}\Omega^{n,0}\},
$$
便得到一个把 $H^*(X;\C)$ 与 Hodge 群
$\bigoplus_{i,j}H^i(\Omega^{j,0})=\bigoplus_{i,j}H^{i,j}(X)$ 联系起来的谱序列
（对 Kähler 流形，这个谱序列著名地退化）。
\end{eg}

希望这次报告已经说明：导出范畴与三角范畴是自然的对象，它们无非是复形构成的
范畴，再把拟同构变成同构；我们可以借拓扑的图景来理解它们（当然略去了一些技术
细节）。遗憾的是，这对真正要在某个 $D(\A)$ 里工作的人来说有点帮倒忙：在那里，
追查一切拟同构、相容性等等是必要的。主要问题在于：虽然锥在三角范畴中只定义到
同构，却无法函子化地选取。这使得许多人相信，还应当有某种更精细的概念有待发展。

\begin{keytakeaways}
\item 导出函子的配方：挑一个适应类 $\RR$，把对象换成 $\RR$-消解，再原样作用
  $F$；所得复形 $\R F(A)$ 在拟同构意义下唯一。
\item 经典导出函子只是“忘了复形”：$R^iF=h^i(\R F)$；长正合列是特异三角的
  上同调影子。
\item 层论中的 Tor、$\R\mathcal Hom$、$\R\Gamma$、超上同调、Leray 与
  local-to-global 谱序列，全是同一台机器的不同输出；de Rham 定理是
  “常数层 $\mathbb R$ 的导出函子”。
\item 收尾的警告：三角范畴里锥只能取到同构，不能函子化；这正是 dg 范畴、
  $A_\infty$ 范畴、稳定 $\infty$-范畴等更精细结构的出发点。
\end{keytakeaways}

\begin{yizhu}[“适应类”的三条各管什么]
第一条（对直和封闭）保证消解可以逐项拼装；第二条（无环进、无环出）是核心：
它保证不同消解给出的 $\R F$ 只差典范同构，即 $\R F$ 不依赖消解；第三条
（每个对象都能嵌入）保证消解存在。对左正合函子，内射对象天然满足这三条；
对右正合函子，投射（或平坦）对象扮演同样的角色。
[来源：Gelfand--Manin III 6.3；Weibel 第 2.3、5.7 节；A 级]
\end{yizhu}

\begin{yizhu}[$\R\Gamma$ 为什么就是层上同调]
$\Gamma$ 是左正合的，内射层构成适应类，故
$\R\Gamma(\F)=\Gamma(I\udot)$，其中 $I\udot$ 是内射消解；取上同调即
$R^i\Gamma(\F)=H^i(X;\F)$。这正是层上同调的标准定义。作者说“$\R\Gamma$ 的
上同调就是标准层上同调”，指的就是这件事。
[来源：Hartshorne《Residues and Duality》II.2；A 级]
\end{yizhu}

\begin{yizhu}[谱序列的读法：$E_2$ 页与收敛]
文中把谱序列写成 $R^iF(R^jG)\Rightarrow R^{i+j}(F\circ G)$，这是“$E_2$ 页”
的简写：$E_2^{i,j}=R^iF(R^jG)$，微分 $d_2$ 走 $(i,j)\to(i+2,j-1)$，最终收敛
到分次对象 $R^{i+j}(F\circ G)$ 的某个滤过。Leray 谱序列对应
$E_2^{i,j}=H^i(Y,R^jp_*)\Rightarrow H^{i+j}(X)$；局部到整体谱序列对应
$E_2^{i,j}=H^i(X,\mathcal Ext^j)\Rightarrow\Ext^{i+j}$。原文此处只写了收敛
箭头的两端，记号略残，见校勘 \ref{kanwu:ss}。
[来源：Weibel 第 5.6 节；Hartshorne《Residues and Duality》II.7；A 级]
\end{yizhu}

\begin{kanwu}[原文把全纯 de Rham 复形称作 Dolbeault 复形]\label{kanwu:dolbeault}
原文的复形
$\C\to\{\OO\xrightarrow{\partial}\Omega^{1,0}\xrightarrow{\partial}\Omega^{2,0}
\to\ldots\to\Omega^{n,0}\}$ 带的是全纯外微分 $\partial$，这是\emph{全纯
de Rham 复形}（holomorphic de Rham complex）；真正的 Dolbeault 复形是
$(\Omega^{0,\bullet},\bar\partial)$，即
$\C\to\{\OO\xrightarrow{\bar\partial}\Omega^{0,1}\xrightarrow{\bar\partial}
\Omega^{0,2}\to\ldots\}$。两者给出不同的谱序列：前者是 Hodge--de Rham 谱序列
$E_1^{p,q}=H^q(X,\Omega^p)\Rightarrow H^{p+q}(X,\C)$，后者是 Dolbeault 谱序列
$E_1^{p,q}=H^q(X,\Omega^p)\Rightarrow H^{p+q}(X,\C)$（在 Kähler 情形两者
同构）。另外，原文把 Hodge 群写成
$H^{i,j}(X)=H^i(\Omega^{j,0})$，与常见约定 $H^{p,q}=H^q(X,\Omega^p)$ 的指标
互换；两种约定都有人用，但同一篇文章内应统一。建议正文改为“全纯 de Rham 复形”，
并把 Hodge 群写成 $H^{i,j}=H^j(\Omega^i)$。详见附录 B 第 4 条。
\end{kanwu}

\begin{kanwu}[“内射层 $\subset$ 拟凝聚层”应为“拟凝聚层中的内射对象”]\label{kanwu:inj}
内射层（作为层范畴的对象）一般并不拟凝聚：例如 $\C$ 上的内射层 $\mathbb Q/\Z$
就不是拟凝聚层。作者想说的其实是“拟凝聚层范畴里的内射对象”，即
$\{$injectives of $\mathrm{QCoh}(X)\}\subset\mathrm{QCoh}(X)$。这在 Noether
概形上确实构成适应类（Hartshorne《Residues and Duality》II.7）。
详见附录 B 第 5 条。
\end{kanwu}

\begin{kanwu}[Leray 与局部到整体谱序列的记号残缺]\label{kanwu:ss}
原文写作 $H^{i+j}_X\Leftarrow H^i_Y(R^jp_*)$ 与
$\Ext^{i+j}\Leftarrow H^i(\mathcal Ext^j)$，两端都缺了 $E_2$ 页的指标。
完整写法应为
$$
E_2^{i,j}=H^i(Y,R^jp_*)\Rightarrow H^{i+j}(X),\qquad
E_2^{i,j}=H^i(X,\mathcal Ext^j)\Rightarrow\Ext^{i+j}(A,B).
$$
详见附录 B 第 6 条。
\end{kanwu}

\begin{kanwu}[对偶公式缺“$A$ 是完美复形”的条件]\label{kanwu:perfect}
原文断言 $\R\mathcal Hom(A,B)\cong B\Ltimes A\ve$，其中
$A\ve=\R\mathcal Hom(A,\OO)$。这个公式的标准形式要求 $A$ 是\term{完美复形}
{perfect complex}（局部可由有限秩自由复形表示）：此时
$\R\mathcal Hom(A,B)\cong\R\mathcal Hom(A,\OO)\Ltimes B$ 由求值映射
$A\ve\Ltimes A\to\OO$ 给出。原文对一般的 $A$、$B$ 断言，条件不成立。反例：取
$X$ 为光滑曲线、$x\in X$、$A=B=k_x$（$x$ 处的摩天层）。$\OO_X$ 无挠，故
$\SheafExt^0(k_x,\OO)=0$，曲线上 $\SheafExt^1(k_x,\OO)=k_x$，于是
$A\ve\cong k_x[1]$；又 $k_x\Ltimes k_x\cong k\oplus k[1]$，故右边
$B\Ltimes A\ve\cong k[1]\oplus k[2]$；而左边
$\R\mathcal Hom(k_x,k_x)\cong k\oplus k[1]$。两者差一个平移，不同构。
建议补上“设 $A$ 为完美复形”。详见附录 B 第 7 条。
\end{kanwu}

\begin{conceptbox}{导出函子：一台机器，五种输出}
把 §\ref{sec:df} 的例子排成一张表，可以看出它们都是“$\R$ + 消解”的同一台机器：
\begin{center}\small
\begin{tabular}{@{}lll@{}}
\toprule
函子 & 消解用什么 & 输出 \\
\midrule
$\Gamma$ & 内射层（或细层） & 层上同调 $H^i(X;\F)$ \\
$\otimes$ & 平坦层／局部自由层 & $\Tor_i$ \\
$\mathcal Hom$ & 局部自由（第一个变量）／内射（第二个） & $\Ext^i$ \\
$\Gamma\circ\mathcal Hom$ & 局部自由 + 内射 & 局部到整体谱序列 \\
$\Gamma$（作用到复形） & 内射消解 & 超上同调 $\mathbb H^i$ \\
\bottomrule
\end{tabular}
\end{center}
作者反复强调的“复形好、上同调坏”，在最后一行的对照下最清楚：超上同调保留了
复形层面的信息，而 $H^i$ 只留下一个分次群。
[来源：原文 §\ref{sec:df} 的例 6.2--6.6；Gelfand--Manin 第 III 章 §3]
\end{conceptbox}
